% Part: model-theory
% Chapter: models-of-arithmetic
% Section: non-standard-models

\documentclass[../../../include/open-logic-section]{subfiles}

\begin{document}

\olfileid{mod}{mar}{nst}
\section{Model Nonstandar}

\begin{explain}
Kita nyebut !!a{structure}
kanggo $\Lang{L_A}$ standar
yen isomorfis
karo~$\Struct{N}$.
Yen !!a{structure}
ora isomorfis
karo~$\Struct{N}$,
struktur iku diarani
nonstandar.
\end{explain}

\begin{defn}
!!{structure}~$\Struct{M}$
kanggo $\Lang{L_A}$
iku \emph{nonstandar}
yen ora isomorfis
karo~$\Struct{N}$.
!!{element}s
$x \in \Domain{M}$
kang padha karo
$\Value{\num{n}}{M}$
kanggo sawenehing
$n \in \Nat$
diarani \emph{wilangan standar}
(saka $\Struct{M}$),
dene kang ora mangkono
diarani \emph{wilangan nonstandar}.
\end{defn}

\begin{explain}
Miturut
\olref[stm]{prop:standard-domain},
saben !!{structure} standar
kanggo~$\Lang{L_A}$
mung ngemot !!{element}s
standar. Mula, yen
!!{structure} ngemot
!!{element} nonstandar,
!!{structure} iku mesthi
nonstandar.
Kosok baline ora mesthi:
struktur nonstandar bisa
mung ngemot unsur standar,
kaya kang dituduhake
dening gladhen ing bagean
sadurunge.
Cathetan koreksi sumber OLPL-097:
sumber mratelakake kosok baline
kang ora bener tanpa
hipotesis tambahan.
\end{explain}

\begin{prop}
Yen !!a{structure}~$\Struct{M}$
kanggo $\Lang{L_A}$
ngemot wilangan nonstandar,
$\Struct{M}$ iku nonstandar.
\end{prop}

\begin{proof}
Upamakna ora mangkono,
yaiku $\Struct{M}$
standar nanging ngemot
wilangan nonstandar~$x$.
Jupuken isomorfisme
$g\colon \Nat \to \Domain{M}$.
Kanthi gampang bisa
dideleng (lumantar induksi
tumrap~$n$) yen
$g(\Value{\num{n}}{N}) =
\Value{\num{n}}{M}$.
Kanthi tembung liya,
$g$ nggawa wilangan
standar saka~$\Struct{N}$
menyang wilangan standar
saka~$\Struct{M}$.
Yen $\Struct{M}$
ngemot wilangan nonstandar,
$g$ ora bisa
!!{surjective},
kosok balen karo hipotesis.
\end{proof}

\begin{prob}
Elinga yen $\Th{Q}$
ngemot aksioma-aksioma
\begin{align*}
& \lforall[x][\lforall[y][(\eq[x'][y'] \lif \eq[x][y])]] \tag{$!Q_1$}\\
& \lforall[x][\eq/[\Obj 0][x']] \tag{$!Q_2$}\\
& \lforall[x][(\eq[x][\Obj 0] \lor \lexists[y][\eq[x][y']])] \tag{$!Q_3$}
\end{align*}
Wenehana !!{structure}s~$\Struct{M_1}$,
$\Struct{M_2}$, $\Struct{M_3}$
saéngga
\begin{enumerate}
\item $\Sat{M_1}{!Q_1}$, $\Sat{M_1}{!Q_2}$, $\Sat/{M_1}{!Q_3}$;
\item $\Sat{M_2}{!Q_1}$, $\Sat/{M_2}{!Q_2}$, $\Sat{M_2}{!Q_3}$; lan
\item $\Sat/{M_3}{!Q_1}$, $\Sat{M_3}{!Q_2}$, $\Sat{M_3}{!Q_3}$;
\end{enumerate}
Cetha yen kanggo saben
struktur iku, sampeyan mung
kudu nemtokake
$\Assign{\Obj{0}}{M_i}$
lan $\Assign{\prime}{M_i}$.
\end{prob}

\begin{explain}
Cukup gampang nemtokake
!!{structure}s nonstandar
kanggo $\Lang{L_A}$.
Upamane, jupuken
struktur kanthi
!!{domain}~$\Int$
lan tafsirna kabeh
simbol nonlogis
kaya lumrahe.
Amarga wilangan negatif
dudu nilai $\num{n}$
kanggo sembarang~$n$,
struktur iki nonstandar.
Mesthi struktur iki
ora dadi \emph{model}
aritmetika ing teges
ndadekake ukara-ukara
kang padha bener
kaya~$\Struct{N}$.
Upamane,
$\lforall[x][\eq/[x'][\Obj{0}]]$
iku luput. Nanging
anane model aritmetika
nonstandar bisa dibuktekake
kanthi gampang
nganggo teorema kompaktitas.
\end{explain}

\begin{prop}
Upamakna
$\Th{TA} = \Setabs{!A}{\Sat{N}{!A}}$
iku teori~$\Struct{N}$.
$\Th{TA}$ nduweni
model nonstandar
kang !!a{enumerable}.
\end{prop}

\begin{proof}
Jembarake $\Lang{L_A}$
nganggo !!{constant} anyar~$c$
lan gatekna himpunan
!!{sentence}s
\[
\Gamma = \Th{TA} \cup \{\eq/[c][\num{0}], \eq/[c][\num{1}],
\eq/[c][\num{2}], \dots\}
\]
Sembarang model~$\Struct{M^c}$
saka~$\Gamma$ bakal ngemot
!!a{element}~$x =
\Assign{c}{M^c}$
kang nonstandar, amarga
$x \neq \Value{\num{n}}{M}$
kanggo saben
$n \in \Nat$.
Kajaba iku, cetha yen
$\Sat{M^c}{\Th{TA}}$,
amarga $\Th{TA} \subseteq \Gamma$.
Yen $\Struct{M^c}$
diowahi dadi !!a{structure}~$\Struct{M}$
kanggo $\Lang{L_A}$
kanthi mung nglalekake~$c$,
ranahé isih ngemot
$x$ nonstandar iku,
lan uga
$\Sat{M}{\Th{TA}}$.
Bab pungkasan iki
dijamin amarga $c$
ora ana ing~$\Th{TA}$.
Mula cukup nuduhake
yen $\Gamma$ nduweni model.
Cathetan koreksi sumber OLPL-098:
interpretasi konstanta anyar
kudu dijupuk saka struktur
kang isih ngemot konstanta iku.

Kita nggunakake
teorema kompaktitas kanggo
nuduhake yen~$\Gamma$
nduweni model.
Yen saben subhimpunan
winates saka~$\Gamma$
bisa dipenuhi,
mula $\Gamma$ uga
bisa dipenuhi.
Gatekna sembarang
subhimpunan winates
$\Gamma_0 \subseteq \Gamma$.
$\Gamma_0$ ngemot
sawenehing !!{sentence}s
saka~$\Th{TA}$
lan sawenehing kang
awujud~$\eq/[c][\num{n}]$,
nanging cacahe mung winates.
Yen ana syarat wujud iku,
upamakna $k$ iku
indeks paling gedhe
saéngga
$\eq/[c][\num{k}] \in \Gamma_0$;
yen ora ana, pilih
indeks nol minangka gantine.
Tetepna $\Struct{N_k}$
kanthi njembarake~$\Struct{N}$
supaya ngemot tafsir
$\Assign{c}{N_k} = k+1$.
$\Sat{N_k}{\Gamma_0}$:
yen $!A \in \Th{TA}$,
mula $\Sat{N_k}{!A}$
amarga $\Struct{N_k}$
padha karo~$\Struct{N}$
ing kabeh bab kajaba~$c$,
lan $c$ ora ana
ing~$!A$.
Lan
$\Sat{N_k}{\eq/[c][\num{n}]}$,
amarga $n \le k$,
lan
$\Value{c}{N_k} = k+1$.
Mula saben subhimpunan
winates saka~$\Gamma$
bisa dipenuhi.
Cathetan koreksi sumber OLPL-099:
yen subhimpunan winates
ora ngemot syarat pengecualian,
indeks paling gedhe
ora ana lan kudu dipilih
kanthi kapisah.
Amarga basa kang ditambahi
konstanta iki bisa didhaptar
lan modelé tanpa wates,
teorema Löwenheim--Skolem mudhun
menehi model kang bisa didhaptar.
Redukté marang basa asal
isih dadi model nonstandar
kang bisa didhaptar.
Cathetan koreksi sumber OLPL-100:
langkah iki nggenepi
pratelan bisa didhaptar
ing proposisi.
\end{proof}

\end{document}
